\documentclass[11pt,a4paper]{article}

\usepackage[T1]{fontenc}
\usepackage[utf8]{inputenc}
\usepackage{lmodern}
\usepackage{amssymb,amsmath,amsthm,mathtools}
\usepackage[margin=2.5cm]{geometry}
\usepackage{parskip}
\usepackage[hidelinks]{hyperref}
\usepackage{microtype}

% Theorem environments (arXiv / journal style)
\theoremstyle{plain}
\newtheorem{theorem}{Theorem}[section]
\newtheorem{proposition}[theorem]{Proposition}
\newtheorem{lemma}[theorem]{Lemma}
\newtheorem{corollary}[theorem]{Corollary}

\theoremstyle{definition}
\newtheorem{definition}[theorem]{Definition}
\newtheorem{example}[theorem]{Example}
\newtheorem{remark}[theorem]{Remark}

% Notation
\newcommand{\Type}{\mathbf{Type}}
\newcommand{\id}{\mathrm{id}}
\newcommand{\Cont}{\mathbf{Cont}}
\newcommand{\Shape}{\mathrm{Shape}}
\newcommand{\Pos}{\mathrm{Pos}}
\newcommand{\Ext}{\mathrm{Ext}}
\newcommand{\seq}{\mathbin{\triangleleft}}
\newcommand{\Dir}{\mathbin{\otimes}}

\title{The Dirichlet Tensor on Containers}
\author{Extracted from a machine-checked development}
\date{}

\begin{document}
\maketitle

\begin{abstract}
We equip the category $\Cont$ of containers with the Dirichlet tensor
\[
  C \Dir D
  \;=\;
  \bigl(S \times T,\; (s,t) \mapsto P(s) \times Q(t)\bigr)
\]
and unit $y$ (one shape, one position---the same container that is the
unit for sequential composition). Bifunctoriality, structural
isomorphisms, naturality, and the Mac Lane coherence laws all hold by
direct computation. The extension of a Dirichlet tensor has a useful
normal form but is \emph{not} a pointwise combination of the factor
extensions; the honest monoidal comparison is Day convolution of
$(\mathbf{Set}, \times, 1)$, with a coend isomorphism (Niu--Spivak). We
also record the uniformity morphism $C \Dir D \to C \seq D$ embedding
the Dirichlet tensor as the non-dependent fragment of sequential
composition, and prove it is not an isomorphism in general. The proofs
of all stated results have been machine-checked; we record only the
mathematical content.
\end{abstract}

\tableofcontents
\bigskip

%======================================================================
\section{Introduction}
%======================================================================

Containers present polynomial endofunctors. Among the monoidal structures
on the category $\Cont$ of containers, two share shapes $S \times T$ and
differ only in positions:
\begin{align}
  C \times D
  &\;=\;
  \bigl(S \times T,\; (s,t) \mapsto P(s) \oplus Q(t)\bigr)
  &&\text{(categorical product)}, \\
  C \Dir D
  &\;=\;
  \bigl(S \times T,\; (s,t) \mapsto P(s) \times Q(t)\bigr)
  &&\text{(Dirichlet tensor)}.
\end{align}
That single swap---sum of positions versus product of positions---is the
heart of the chapter. On extensions it is the difference between a
pointwise product of functors (by the exponential law
$X^{A+B} \cong X^A \times X^B$) and a structure that does \emph{not}
split pointwise (witness: $y^3 \Dir y^3 = y^9$, while anything pointwise
would give $y^6$).

The unit of the Dirichlet tensor is $y$, the one-shape one-position
container---definitionally the same as the unit~$I$ of sequential
composition~$\seq$. Two of the four monoidal structures on $\Cont$ thus
share a unit; both restrict to the identity on that container.

\paragraph{What this note claims.}
\begin{enumerate}
  \item $(\Cont, \Dir, y)$ is a monoidal category, with all laws
        transport-free.
  \item The extension of $C \Dir D$ has the normal form
        $\sum_{(s,t)} X^{P(s)\times Q(t)}$, a definitional equality of
        types and \emph{not} a strict-monoidality statement.
  \item The uniformity morphism $C \Dir D \to C \seq D$ is canonical but
        not an isomorphism in general.
\end{enumerate}

We work throughout in a single universe of types. Proofs are indicated by
the phrase ``by lean.''

%======================================================================
\section{The tensor and its unit}
%======================================================================

\begin{definition}[Container]\label{def:container}
A \textbf{container} $C = (S,P)$ consists of a shape type~$S$ and a
position family~$P$. Morphisms are shape maps with contravariant position
maps; we write $\Cont$ for the category.
\end{definition}

\begin{definition}[Dirichlet tensor]\label{def:dirichlet}
The \textbf{Dirichlet tensor} of containers $C = (S,P)$ and $D = (T,Q)$
is the container
\begin{equation}
  C \Dir D
  \;=\;
  \bigl(
    S \times T,\;
    (s,t) \mapsto P(s) \times Q(t)
  \bigr).
\end{equation}
Positions are a \emph{product}---this is exactly what distinguishes
$\Dir$ from the categorical product, whose positions are the
\emph{coproduct} $P(s) \oplus Q(t)$.
\end{definition}

\begin{definition}[Unit]\label{def:y}
The unit for the Dirichlet tensor is the container
\begin{equation}
  y \;=\; \bigl(1,\; {\ast} \mapsto 1\bigr)
\end{equation}
---one shape, one position. It is the same container as the unit~$I$ of
the sequential operator~$\seq$; the name~$y$ is the polynomial reading
as the representable $y^1$.
\end{definition}

\begin{remark}
The left and right unitors for $\Dir$ are unit-collapses in both shape
and position components: $y \Dir C$ has shapes $1 \times S$ and
positions $1 \times P(s)$.
\end{remark}

%======================================================================
\section{Bifunctoriality}
%======================================================================

\begin{definition}[Action on morphisms]\label{def:dir2}
Given morphisms $\mu : C \to C'$ and $\nu : D \to D'$, the
\textbf{Dirichlet action}
\[
  \mu \Dir_2 \nu
  \;:\;
  C \Dir D \;\longrightarrow\; C' \Dir D'
\]
acts componentwise:
\begin{align}
  \text{on shapes:}\quad
  (s,t) &\;\longmapsto\; \bigl(\mu(s),\, \nu(t)\bigr), \\
  \text{on positions:}\quad
  (p,q) &\;\longmapsto\;
  \bigl(\mu^\sharp_s(p),\, \nu^\sharp_t(q)\bigr).
\end{align}
Unlike sequential composition, whose morphism action threads the outer
position map through the inner shape assignment, the Dirichlet action is
strictly componentwise.
\end{definition}

\begin{proposition}[Functor laws]\label{prop:dir2-laws}
\begin{enumerate}
  \item $\id_C \Dir_2 \id_D = \id_{C \Dir D}$;
  \item $(\mu \mathbin{;}\mu') \Dir_2 (\nu \mathbin{;}\nu')
        = (\mu \Dir_2 \nu) \mathbin{;} (\mu' \Dir_2 \nu')$.
\end{enumerate}
\end{proposition}

\begin{proof}
By lean.
\end{proof}

%======================================================================
\section{Structural isomorphisms and coherence}
%======================================================================

\begin{definition}[Associator and unitors]\label{def:dir-struct}
\begin{enumerate}
  \item The \textbf{associator}
        $(C \Dir D) \Dir E \xrightarrow{\sim} C \Dir (D \Dir E)$ is
        product reassociation on shapes and (backward) product
        reassociation on positions. Product reassociation is its own
        mirror image, so the position component is transport-free.
  \item The \textbf{left unitor} $y \Dir C \xrightarrow{\sim} C$ and
        \textbf{right unitor} $C \Dir y \xrightarrow{\sim} C$ are the
        unit-collapses in both components.
\end{enumerate}
Both round trips of each isomorphism hold by definitional $\eta$ for
products and the unit type.
\end{definition}

\begin{proposition}[Naturality]\label{prop:dir-nat}
The associator and both unitors of the Dirichlet tensor are natural in
all arguments.
\end{proposition}

\begin{proof}
By lean.
\end{proof}

\begin{theorem}[Coherence]\label{thm:dir-coherence}
The pentagon and triangle identities hold for the Dirichlet tensor.
\end{theorem}

\begin{proof}
By lean.
\end{proof}

\begin{remark}
Every map in sight is a product/unit shuffle. Definitional $\eta$ for
structures identifies the two legs of each coherence diagram on the nose.
\end{remark}

\begin{theorem}[$(\Cont, \Dir, y)$ is monoidal]\label{thm:dir-monoidal}
The category of containers, equipped with the Dirichlet tensor and unit
$y$, is a monoidal category.
\end{theorem}

\begin{proof}
By lean.
\end{proof}

\begin{remark}
As with the coproduct and product, a formal development that admits only
one canonical monoidal instance on~$\Cont$ packages the Dirichlet
structure as named data rather than a global instance. The sequential
operator is preferred as the canonical instance because its comonoids are
directed containers. Mathematically, four monoidal structures on one
category is the point.
\end{remark}

%======================================================================
\section{Extension: a normal form, not a strictness theorem}
%======================================================================

\begin{theorem}[Normal form]\label{thm:normal-form}
For all containers $C = (S,P)$, $D = (T,Q)$ and every type~$X$,
\begin{equation}
  \Ext(C \Dir D)\,X
  \;=\;
  \sum_{(s,t) \,:\, S \times T}
  \bigl(P(s) \times Q(t) \to X\bigr)
\end{equation}
as a definitional equality of types.
\end{theorem}

\begin{proof}
By lean.
\end{proof}

\begin{remark}[What this is not]
Theorem~\ref{thm:normal-form} is a useful computation rule and nothing
more. It is \emph{not} a monoidal-functor statement, and in particular it
does not say that $\Ext(-)$ is strict monoidal for~$\Dir$.

A comparison map must relate $\Ext(C \Dir D)$ to $\Ext\,C$ and $\Ext\,D$
through a tensor on the target category. The right-hand side above is not
built from the functors $\Ext\,C$ and $\Ext\,D$ at all: it re-reads the
container's presentation $(S,T,P,Q)$. Observing that $\Ext(-)$ ``preserves''
a tensor that was itself defined by reading the presentation is circular;
the equality is an unfolding of a definition.
\end{remark}

\begin{remark}[Honest comparison: Day convolution]
Compare the three comparisons with real content:
\begin{itemize}
  \item sequential:
        $\Ext(G \seq F)\,X \cong \Ext\,G\bigl(\Ext\,F\,X\bigr)$
        (target tensor $=$ composition of functors; a genuine iso);
  \item coproduct:
        $\Ext(C+D)\,X \cong \Ext\,C\,X \oplus \Ext\,D\,X$
        (target tensor $=$ pointwise coproduct);
  \item product:
        $\Ext(C \times D)\,X \cong \Ext\,C\,X \times \Ext\,D\,X$
        (target tensor $=$ pointwise product).
\end{itemize}
The correct intrinsic account of~$\Dir$ is due to Niu and Spivak:
$\Dir$ is the \textbf{Day convolution} on $[\mathbf{Set},\mathbf{Set}]$ of
the cartesian monoidal structure $(\mathbf{Set}, \times, 1)$. The
comparison map $\Ext\,C \,\odot\, \Ext\,D \to \Ext(C \Dir D)$ is the
co-Yoneda/density isomorphism out of a coend---an isomorphism, but
emphatically not an identity. Thus $\Ext(-)$ is \textbf{strong}, not
strict, monoidal for~$\Dir$.
\end{remark}

\begin{remark}[Why $\Dir$ is non-pointwise]
Being a Day convolution is \emph{not} what makes $\Dir$ non-pointwise.
The categorical product~$\times$ on $\Cont$ is also a Day convolution---of
$(\mathbf{Set}, +, 0)$---and it \emph{is} pointwise. The real criterion is
whether the corepresentable splits the domain tensor:
\begin{itemize}
  \item $y^{A+B} \cong y^A \times y^B$ --- so Day-of-$+$ (the product
        $\times$) is pointwise;
  \item $y^{A \times B}$ admits \textbf{no} such splitting --- so Day-of-
        $\times$ (the Dirichlet $\Dir$) is not.
\end{itemize}
Witness: $y^3 \Dir y^3 = y^9$, whereas anything pointwise would give
$y^6$.
\end{remark}

\paragraph{Taxonomy of the four monoidal structures on $\Cont$.}
\begin{center}
\begin{tabular}{llll}
structure & Day convolution of & unit & $\Ext$ pointwise? \\
\hline
product $\times$ & $(\mathbf{Set}, +, 0)$ & $1 = y^0$ & yes \\
Dirichlet $\Dir$ & $(\mathbf{Set}, \times, 1)$ & $y = y^1$ & no \\
coproduct $+$ & not a Day convolution & $0$ & yes \\
sequential $\seq$ & not a Day convolution & $I = y$ & no
\end{tabular}
\end{center}

%======================================================================
\section{Dirichlet as the uniform fragment of sequential composition}
%======================================================================

Both $C \Dir D$ and $C \seq D$ ``put a $D$ at every $C$-position'', but
\begin{itemize}
  \item $C \seq D$ lets the inner $D$-shape \emph{depend} on the outer
        $C$-position: its shapes are $\sum_{s \,:\, S} (P(s) \to T)$;
  \item $C \Dir D$ forces one inner shape \emph{uniformly}: its shapes
        are $S \times T$.
\end{itemize}
Choosing the constant assignment $p \mapsto t$ therefore embeds $\Dir$
into~$\seq$. On positions the two agree definitionally:
\begin{equation}
  \sum_{p \,:\, P(s)} Q\bigl(({\ast} \mapsto t)(p)\bigr)
  \;=\;
  P(s) \times Q(t).
\end{equation}

\begin{definition}[Uniformity morphism]\label{def:dir-to-seq}
The canonical \textbf{uniformity morphism}
\[
  C \Dir D \;\longrightarrow\; C \seq D
\]
sends the shape $(s,t)$ to $\bigl\langle s,\; {\ast} \mapsto t \bigr\rangle$
(the constant inner-shape assignment); on positions it is the identity,
the fibres agreeing definitionally.
\end{definition}

\begin{theorem}[Uniformity is not an isomorphism]\label{thm:not-iso}
The uniformity morphism is not an isomorphism in general. Concretely,
for $C = y^2$ (one shape, two positions) and $D = 2$ (two shapes, no
positions), the shape map
\[
  2 \;\longrightarrow\; (2 \to 2),
  \qquad
  t \;\longmapsto\; ({\ast} \mapsto t)
\]
misses the identity function on the two-element set, hence is not
surjective on shapes. Equivalently: $y^2 \Dir 2$ has two shapes while
$y^2 \seq 2$ has four; both have no positions, so their extensions at
any~$X$ have two and four elements respectively.
\end{theorem}

\begin{proof}
By lean.
\end{proof}

\begin{remark}
Thus $\Dir \neq \seq$: the uniformity morphism is a genuine comparison,
not an equivalence. The Dirichlet tensor is sequential composition with
the dependency switched off---the uniform (non-dependent) fragment
of~$\seq$.
\end{remark}

%======================================================================
\section{Discussion}
%======================================================================

\paragraph{What has been established.}
The Dirichlet tensor is a monoidal structure on~$\Cont$, transport-free
and coherent. Its extension has a clean normal form that must not be
misread as strict monoidality. The honest comparison is Day convolution.
The uniformity morphism relates $\Dir$ to~$\seq$ without collapsing them.

\paragraph{Scope.}
The coend comparison of Niu--Spivak is not formalised here (it requires
quotients beyond a pure core type theory). The taxonomy of the four
structures is stated for orientation; the coproduct and product monoidal
data are in a companion note, and the sequential structure in another.

\paragraph{Machine checking.}
Every definition and theorem statement above corresponds to a declaration
in a machine-checked formalisation. The phrase ``by lean'' indicates that
the proof has been verified by a proof assistant.

%======================================================================
\section*{Acknowledgements}
%======================================================================

The mathematical content follows the theory of containers of Abbott,
Altenkirch, and Ghani, and the Dirichlet/Day-convolution account of
polynomial functors due to Niu and Spivak. The presentation here is a
prose rendering of a formal development intended for human review of the
theorem statements.

\begin{thebibliography}{9}

\bibitem{abbott2005}
M.~Abbott, T.~Altenkirch, and N.~Ghani.
\newblock Containers: constructing strictly positive types.
\newblock \emph{Theoretical Computer Science}, 342(1):3--27, 2005.

\bibitem{abbott2003}
M.~Abbott, T.~Altenkirch, and N.~Ghani.
\newblock Categories of containers.
\newblock In \emph{FoSSaCS}, volume~2620 of \emph{LNCS}, pages 23--38.
\newblock Springer, 2003.

\bibitem{niu-spivak2023}
N.~Niu and D.~I.~Spivak.
\newblock Polynomial functors: a mathematical theory of interaction.
\newblock arXiv:2312.00990, 2023.

\bibitem{maclane1971}
S.~Mac~Lane.
\newblock \emph{Categories for the Working Mathematician}.
\newblock Springer, 1971.

\end{thebibliography}

\end{document}
